\par\clearpage
\noindent\makebox[0pt][l]{\fontsize{13}{16}\selectfont 416}\hfill{\small\textbf{NACHLASS.}}\hfill\null\par\medskip
\noindent dings wichtig sein, die südlichen grossen Dreiecke noch in mehrere kleinere
zu zerlegen, und durch Messungen an neuen eingeschalteten Standpunkten
sichere und zureichende Grundlagen für jene Aufnahmen zu gewinnen. Auf
die Kosten \textls[120]{dieser} Operationen habe ich bei der obigen Schätzung keine
Rücksicht nehmen können: ihre Veranschlagung würde fast noch misslicher,
aber auf jeden Fall können sie doch, vergleichungsweise gegen die Kosten
neuer grosser Triangulirungen in den noch nicht berührten Landestheilen, nur
klein sein.

Dass es übrigens in Zukunft wünschenswerth sein wird, die Resultate der
Lage aller scharf bestimmten Punkte, wenn sie erst ein geschlossenes Ganzes
bilden, \textls[120]{öffentlich} bekannt zu machen, brauche ich nicht zu bemerken. Von
dem eigentlich rein wissenschaftlichen Theile der bisherigen Messungen versteht
sich dies ohnehin von selbst.\quad\mbox{.\ .\ .\ .\ .\ .}

\bigskip
\begin{center}
[Aus dem Bericht vom 26.~Junius 1828, »die Fortsetzung der Gradmessungs-\par
arbeiten betreffend.«]
\end{center}

\par\noindent\hspace*{\parindent}\mbox{.\ .\ .\ .\ .\ .}\quad Da die Detailaufnahme der noch nicht vermessenen Landestheile
auf die trigonometrischen Operationen gegründet werden soll, und beide Geschäfte
rücksichtlich der zu verwendenden Geldmittel von einander abhängig
sein werden, so war zuvörderst eine ungefähre Überschlagung der Gesammtkosten
erforderlich. Nach einer mir vom Herrn Geh.~Cabinetsrath \textsc{Hoppenstedt}
mitgetheilten Notiz würde der Flächeninhalt der im Detail aufzunehmenden
Landestheile etwa 144 Quadratmeilen betragen; die Kosten der
Detailaufnahme durch Generalstabsofficiere werden auf 200---250 Thaler für
jede Quadratmeile geschätzt, wozu noch etwa 21 Thaler wegen der Copirungskosten
der Karte in 4 Exemplaren hinzu zu rechnen sein würden. Würden also
zusammen 250 Thaler auf die Quadratmeile gerechnet, so würden diese Kosten
etwa 36000 Thaler, folglich mit Inbegriff der Triangulirungskosten in den
von der Gradmessung noch nicht berührten Landestheilen gegen 50000 Thaler
betragen. Es möchten dazu noch ein oder ein paar tausend Thaler zu rechnen
sein wegen der Operationen, die erforderlich sein werden, um innerhalb der
grossen südlichen Dreiecke der Gradmessung eine hinlänglich grosse Anzahl
\par
\clearpage
\noindent\null\hfill{\small\textbf{ZUR HANNOVERSCHEN TRIANGULATION.}}\hfill\makebox[0pt][r]{\fontsize{13}{16}\selectfont 417}\par\medskip
\noindent fester Punkte für die Detailaufnahme festzulegen, worüber ich mich bereits
früher in der im November vorigen Jahres eingereichten Eingabe ausführlicher
erklärt habe. Es würde daher, wenn zu diesen Geschäften jährlich wirklich
5000 Thaler verwendet werden können (was ausser den Geldmitteln auch
von der steten Disponibilität des Personals abhängen wird), zur völligen Vollendung
ungefähr ein Zeitraum von 10 Jahren erforderlich sein.

Nach einem von mir selbst früher gemachten, obwohl vielleicht minder
zuverlässigen Überschlage wäre der Flächeninhalt der im Detail aufzunehmenden
Landestheile 173 Quadratmeilen; danach würde der Überschlag für den
Kostenaufwand etwa 7000 Thaler grösser und die Zeitdauer 1 bis 2 Jahr
länger ausfallen. Und so würde auch, wenn die bestimmten 5000 Thaler
nicht als alljährlich wirklich oder im Durchschnitt zu verwenden, sondern nur
als jedesmaliges Maximum zu betrachten sind, eine verhältnissmässige Verlängerung
der Dauer des ganzen Geschäfts die Folge sein.

Was die Eintheilung der trigonometrischen Arbeiten auf die einzelnen
Jahre betrifft, so möchte es, um diese als Grundlage der Detailaufnahme
schneller vollenden zu können, rathsam sein, anfangs den grössern Theil der
Geldmittel auf dieselbe und den kleinern auf die Detailaufnahme, vielleicht
in dem Verhältniss von \(\frac{2}{3}\) zu \(\frac{1}{3}\), zu verwenden: letztere würde dann in den
spätern Jahren, wo überall eine sichere Grundlage vorhanden ist, wo die Arbeiter
nach und nach immer mehr eingeübt sind, und wo die Geldmittel
allein darauf verwendet werden können, eines um so raschern Fortschreitens
gewiss sein. In Beziehung auf die Anordnung der Reihenfolge der trigonometrischen
Arbeiten ist, wie die Sachen gegenwärtig stehen, weiter kein
Grund vorhanden, eine der andern vorzuziehen, als dass nur darauf gesehen
werden muss, dass die gleichzeitige Detailaufnahme stets mit hinreichendem
Stoff an zuverlässig bestimmten Punkten versehen sei, damit dieselbe niemals
in Gefahr komme, aus Mangel an solchem in Stocken zu gerathen. Eine
speciellere Bestimmung möchte wohl \textls[120]{für den Augenblick}, theils unthunlich,
theils unnöthig, theils nicht einmal rathsam sein, weil ein gewisser Grad
von Freiheit, das den jedesmaligen Umständen nach zweckmässigste zu bearbeiten,
dem schnellern und bessern Fortschreiten nur förderlich sein kann.
\par\noindent\mbox{.\ .\ .\ .\ .\ .}\par
\medskip\noindent\hspace*{0.03\linewidth}{\footnotesize IX.}\hfill\makebox[0.18\linewidth][l]{53}\par
